\begin{lemma}[Scheide maximal-tree package (e)]
Let $(G,e_0,\varphi)$ be a $k$-triple.  Let $T$ be fixed-pair maximal for
$(e_0,\varphi)$, and let
$S=(y_0,e_1,y_1,\ldots,e_n,y_n)$ be any Tashkinov tree for the same pair.
If $i<j$, $\gamma$ is missing at $y_i$, and $\delta$ is missing at $y_j$,
then $\gamma\ne\delta$ and the $(\gamma,\delta)$-component through $y_i$
is an even multigraph path with ends $y_i,y_j$ whose vertices lie in
$V(T)$.  In the bipartite incidence graph this is a simple spanning path of
length divisible by four.  If $\gamma$ is unused through $y_j$, switching
this component preserves a proper coloring with exactly $e_0$ uncolored and
leaves $S$ a Tashkinov tree.
\end{lemma}
\begin{proof}
Put $X=V(T)$ and $X'=V(S)$.  The properness and unique-uncolored-edge
conditions are supplied by \noderef{KTriple}, and missing colors have the
meaning in \noderef{MissingColors}.  By \noderef{TashTree}, $X'$ is elementary.
Since $y_i\ne y_j$ by vertex injectivity in \noderef{TreeSequence}, the
colors missing at those two vertices are distinct; thus
$\gamma\ne\delta$.  Also \noderef{ScheideMaximalContainsEveryTree} gives
$X'\subseteq X$, and \noderef{ScheideMaximalElementaryClosed} says that
$X$ is elementary and closed.

Consider the $(\gamma,\delta)$-component of \noderef{TwoColorComponent}
containing $y_i$ in the bipartite incidence graph
\noderef{TwoColorIncidenceGraph}.  Because $\gamma$ is
missing at $y_i$ and $X$ is elementary, $\delta$ is present at $y_i$ and
properness gives exactly one incident edge of that color.  Thus $y_i$ has
degree one in the two-color component, and its incident edge has color
$\delta$.  Symmetrically, $y_j$ has degree one and its incident edge has
color $\gamma$.  At every other vertex of $X$, both colors are present:
each is missing somewhere in $X$, so elementariness forbids it from being
missing at another vertex.  Properness gives exactly one incident edge of
each color there.

Closedness of $X$ forbids any $\gamma$- or $\delta$-colored boundary edge,
because both colors are missing on $X$.  Hence the entire component through
$y_i$ lies over vertices of $X$.  By
\noderef{TwoColorComponentSpanningPath}, the finite component has a simple
spanning incidence path.  The degree calculation shows that its only
degree-one vertex-nodes are $y_i,y_j$, so these are the two ends.  Along the
corresponding multigraph path, colors alternate; the first edge at $y_i$
has color $\delta$ and the last edge at $y_j$ has color $\gamma$.  Since
these colors differ, the multigraph path has an even number of edges.  Each
multigraph edge contributes two successive incidences, so the incidence
path length is divisible by four.  This proves all path assertions.

Now suppose $\gamma$ is unused through $y_j$ in the sense of
\noderef{TreeColorUnusedThrough}.  The color $\delta$ is unused there as
well.  Otherwise a colored tree edge before or at $y_j$ with color
$\delta$ would, by \noderef{TashkinovTree}, have a still earlier vertex
missing $\delta$; together with the missing occurrence at $y_j$, this would
contradict the elementariness of $X'$.

Switch $\gamma$ and $\delta$ on the displayed component using
\noderef{TwoColorComponentSwap}.  Properness is preserved by
\noderef{TwoColorComponentSwapProper}.  The uncolored edge $e_0$ is not an
edge-node of the component, because every component edge is colored
$\gamma$ or $\delta$; hence it remains uncolored, and every other edge
remains colored.  No colored tree edge
through $y_j$ changes, since both colors are unused there.  At the path
ends, the missing colors are exchanged: $\delta$ becomes missing at $y_i$
and $\gamma$ at $y_j$; at every other tree vertex the two incident colors
are merely exchanged and its missing-color set is unchanged.  For every
tree edge later than $y_j$, a witness in a color outside
$\{\gamma,\delta\}$ is unchanged.  If its color or its old witness uses
$\gamma$ or $\delta$, the appropriate one of $y_i,y_j$ supplies the new
or unchanged missing-color witness, and both vertices precede that later
edge because $i<j$.  Thus every later tree edge still has an earlier
missing-color witness.  The switched coloring therefore leaves the same
displayed sequence $S$ a Tashkinov tree.
\end{proof}
